%\vspace*{-0.4cm}
\chapter[Thème 2 $-$ Primitives et équa diff]{Thème 2 : Modèles continus \\Primitives et équations différentielles}
\thispagestyle{fancy}

\vspace*{-0.6cm}
\section{Primitive d'une fonction}
\subsection{Définition et propriétés}

\begin{bexemple}
	On considère les fonctions $f$ et $F$ définies par : $f(x)=2 x+3$ et $F(x)=x^{2}+3 x-1$\\
	Si on dérive $F$, on constate que : $F^{\prime}(x)=2 x+3=f(x)$.\\
	Lorsque $F^{\prime}=f$, on dit que $F$ est une primitive de $f$.
\end{bexemple}

\begin{bdefi}
	$f$ est une fonction continue sur un intervalle $I$.\\
	\phantom{$\dfrac{1}{2}$}\pointilles[\textwidth-0.5cm]
\end{bdefi}	

\begin{brmq}
	Dans ces conditions, dire que \guillemets{$F$ est une primitive de $f$} revient à dire que  \guillemets{$f$ est la dérivée de $F$}.
\end{brmq}

\begin{bmethodet}{Vérifier qu'une fonction est une primitive d'une autre fonction}
	\grasInMet{Enoncé : }
	
	Dans chaque cas, dire si $F$ est une primitive de $f$.
		\begin{enumMet}
		\item $F(x)=\dfrac{x^{2}}{2}$ et $f(x)=x$.
		\item $F(x)=x e^{x}$ et $f(x)=e^{x}(x+1)$
		\item $F(x)=\dfrac{\ln (x)}{x}$ et $f(x)=\dfrac{-\ln (x)}{x^{2}}$.
	\end{enumMet}
\end{bmethodet}

\begin{bprop}
	\phantom{$\dfrac{1}{2}$}\pointilles[\textwidth-0.5cm]
	
	\phantom{$\dfrac{1}{2}$}\pointilles[\textwidth-0.5cm]
\end{bprop}

\begin{bdemo}
\end{bdemo}

\begin{bprop}
	$f$ est une fonction continue sur un intervalle $I$.\\
	\phantom{$\dfrac{1}{2}$}\pointilles[\textwidth-0.5cm]
	
	\phantom{$\dfrac{1}{2}$}\pointilles[\textwidth-0.5cm]
\end{bprop}


\begin{bdemo}
\end{bdemo}

\begin{bexemple}
	On a vu dans la méthode précédente que $F$ est une primitive de $f$ avec :\\
	$F(x)=\dfrac{x^{2}}{2}$ et $f(x)=x$.\\
	Donc, la fonction $G$ définie par $G(x)=\dfrac{x^{2}}{2}+5$ est également une primitive de $f$.\\
	En effet : $G^{\prime}(x)=\dfrac{2 x}{2}+0=x=f(x)$.
\end{bexemple}

\begin{bprop}
	\phantom{$\dfrac{1}{2}$}\pointilles[\textwidth-0.5cm]
	
	\phantom{$\dfrac{1}{2}$}\pointilles[\textwidth-0.5cm]
\end{bprop}

\begin{brmq}
	Bien que l'existence étant assurée, la forme explicite d'une primitive n'est pas toujours connue. Par exemple, la fonction $x \mapsto e^{-x^{2}}$ ne possède pas de primitive sous forme explicite.
\end{brmq}

\begin{bmethodet}{Recherche d'une primitive particulière}
	\grasInMet{Enoncé : }
	
	Soit la fonction $f$ définie sur $\mathbb{R}^{*}$ par $f(x)=\dfrac{e^{2 x}(2 x-1)}{x^{2}}$.
	\begin{enumMet}
		\item Démontrer que la fonction $F$ définie sur $\mathbb{R}^{*}$ par $F(x)=\dfrac{e^{2 x}}{x}$ est une primitive de $f$.\\
		\item Déterminer la primitive $G$ de la fonction $f$ qui s' annule en $x=1$.
	\end{enumMet}
	
\end{bmethodet}

\subsection{Primitives des fonctions usuelles}

\begin{center}
	\renewcommand{\arraystretch}{2}
	\begin{tabular}{|Xc|Xc|}
		\hline
		Fonction $\boldsymbol{f}$ & Une primitive $\boldsymbol{F}$ \\
		\hline
		$a, a \in \mathbb{R}$ & \\
		\hline
		\begin{tabular}{c}
			$x^{n}$ \\
			avec $n \in \mathbb{N}^{*}$ \\
		\end{tabular} &  \\
		\hline
		$\dfrac{1}{x^{2}}$ &  \\
		\hline
		$\dfrac{1}{x}$ &  \\
		$\operatorname{Avec} x>0$ &  \\
		\hline
		$\dfrac{1}{\sqrt{x}}$ &  \\
		\hline
		$e^{x}$ &  \\
		\hline
	\end{tabular}
\end{center}

\subsection{Linéarité des primitives}

\begin{bprop}
	Si $F$ est une primitive de $f$ et $G$ est une primitive de $g$ alors :
	\begin{itemProp}
		\item \phantom{$\dfrac{1}{2}$}\pointilles[\textwidth-0.8cm]
		\item \phantom{$\dfrac{1}{2}$}\pointilles[\textwidth-0.8cm]
	\end{itemProp}
\end{bprop}

\begin{bdemo}
	\begin{itemDem}
		\item \phantom{$\dfrac{1}{2}$}\pointilles[\textwidth-0.5cm]
		\item \phantom{$\dfrac{1}{2}$}\pointilles[\textwidth-0.5cm]
	\end{itemDem}
\end{bdemo}

\begin{bmethodet}{Déterminer une primitive (1)}
	\grasInMet{Enoncé : }
	
	Dans chaque cas, déterminer une primitive $F$ de la fonction $f$.
	\begin{multicols}{2}
		\begin{enumMet}
		\item $f(x)=x^{3}-2$
		\item $f(x)=3 x^{2}-\dfrac{3}{x^{2}}$
		\item $f(x)=\dfrac{3}{x}$ sur $] 0 ;+\infty[$
		\item $f(x)=\dfrac{2}{\sqrt{x}}$
	\end{enumMet}
	\end{multicols}
\end{bmethodet}

\subsection{Primitives de fonctions composées}
$u$ est une fonction dérivable sur un intervalle $I$.

\begin{center}
	\renewcommand{\arraystretch}{1.7}
	\begin{tabular}{|Xc|Xc|}
		\hline
		Fonction & Une primitive \\
		\hline
		$2 u^{\prime} u$ & \\
		\hline
		$u^{\prime} e^{u}$ & \\
		\hline
		$\dfrac{u^{\prime}}{u}$ & \\
		avec $u>0$ &  \\
		\hline
	\end{tabular}
\end{center}


\begin{bmethodet}{Déterminer une primitive (2)}
	\grasInMet{Enoncé : }
	
	Dans chaque cas, déterminer une primitive $F$ de la fonction $f$.
		\begin{enumMet}
		\item $f(x)=(2 x-5)\left(x^{2}-5 x+4\right)$
		\item $f(x)=x e^{x^{2}}$
		\item $f(x)=\dfrac{x^{2}}{x^{3}+1}$
	\end{enumMet}
\end{bmethodet}

\section{Équations différentielles}
\subsection{Définition d'une équation différentielle}

\begin{bdefi}
	\phantom{$\dfrac{1}{2}$}\pointilles[\textwidth-0.5cm]
	
	\phantom{$\dfrac{1}{2}$}\pointilles[\textwidth-0.5cm]
\end{bdefi}	

\begin{bexemple}
	\begin{itemDef}
		\item L'équation $f^{\prime}(x)=5$ est une équation différentielle.
		
		L'inconnue est la fonction $f$.\\
		En considérant que $y$ est la fonction inconnue qui dépend de $x$, l'équation peut se noter : $y^{\prime}=5$
		\item L'équation $y^{\prime}=2 x^{2}-3$ est également une équation différentielle.
		L'inconnue est la fonction $y$ dont la dérivée est égale à $2 x^{2}-3$.
	\end{itemDef}	
\end{bexemple}

\subsection{Équation différentielle du type $y^{\prime}=f$}

\begin{bdefi}
	Soit une fonction $f$ définie sur un intervalle $I$.\\
	\phantom{$\dfrac{1}{2}$}\pointilles[\textwidth-0.5cm]
	
	\phantom{$\dfrac{1}{2}$}\pointilles[\textwidth-0.5cm]
\end{bdefi}

\begin{bprop}
	\phantom{$\dfrac{1}{2}$}\pointilles[\textwidth-0.5cm]
	
	\phantom{$\dfrac{1}{2}$}\pointilles[\textwidth-0.5cm]
	
	\phantom{$\dfrac{1}{2}$}\pointilles[\textwidth-0.5cm]
\end{bprop}

\begin{bmethodet}{Vérifier qu'une fonction est solution d'une équation différentielle}
	\grasInMet{Enoncé : }
	
	Prouver que la fonction $g$ définie sur $] 0 ;+\infty\left[\operatorname{par} g(x)=3 x^{2}+\ln x\right.$ est solution de l'équation différentielle $y^{\prime}=6 x+\dfrac{1}{x}$.
\end{bmethodet}

\subsection{Équations différentielles du type $y^{\prime}=a y$}

\begin{bprop}
	\phantom{$\dfrac{1}{2}$}\pointilles[\textwidth-0.5cm]
	
	\phantom{$\dfrac{1}{2}$}\pointilles[\textwidth-0.5cm]
\end{bprop}

\begin{bdemo}
\end{bdemo}

\begin{bmethodet}{Résoudre une équation différentielle du type $y^{\prime}=a y$}
	\grasInMet{Enoncé : }
	
	On considère l'équation différentielle $3 y^{\prime}+5 y=0$.
	
	\begin{enumMet}
		\item \begin{enumMet}
			\item Déterminer la forme générale des fonctions solutions de l'équation.
			\item Représenter à l'aide de la calculatrice ou d'un logiciel, quelques courbes des fonctions solutions.
		\end{enumMet}
		\item Déterminer l'unique solution $f$ telle que $f(1)=2$.
	\end{enumMet}
	
\end{bmethodet}

\begin{bprop}
	\phantom{$\dfrac{1}{2}$}\pointilles[\textwidth-0.5cm]
	
	\phantom{$\dfrac{1}{2}$}\pointilles[\textwidth-0.5cm]
\end{bprop}

\begin{bdemo}
	\begin{itemDem}
		\item \phantom{$\dfrac{1}{2}$}\pointilles[\textwidth-0.5cm]
		\item \phantom{$\dfrac{1}{2}$}\pointilles[\textwidth-0.5cm]
	\end{itemDem}
\end{bdemo}

\subsection{Équations différentielles du type $y^{\prime}=a y+b$}

\begin{bprop}
\phantom{$\dfrac{1}{2}$}\pointilles[\textwidth-0.5cm]

\phantom{$\dfrac{1}{2}$}\pointilles[\textwidth-0.5cm]
\end{bprop}

\begin{bdemo}
\end{bdemo}

\begin{bprop}
	\phantom{$\dfrac{1}{2}$}\pointilles[\textwidth-0.5cm]
	
	\phantom{$\dfrac{1}{2}$}\pointilles[\textwidth-0.5cm]
	
	\phantom{$\dfrac{1}{2}$}\pointilles[\textwidth-0.5cm]
	
	\phantom{$\dfrac{1}{2}$}\pointilles[\textwidth-0.5cm]
\end{bprop}

\begin{brmq}
	L'équation $y^{\prime}=a y+b$ est appelée équation différentielle linéaire du premier ordre à coefficients constants.
\end{brmq}
\begin{bmethodet}{Résoudre une équation différentielle du type $y^{\prime}=a y+b$}
	\grasInMet{Enoncé : }
	
	On considère l'équation différentielle $(E): 2 y^{\prime}-y=3$.
	\begin{enumMet}
		\item Déterminer une solution particulière constante de l'équation différentielle ( $E$ ).\\
		\item Déterminer la forme générale des solutions de l'équation différentielle $y^{\prime}=\dfrac{1}{2} y$.\\
		\item En déduire la forme générale des solutions de l'équation différentielle ( $E$ ).\\
		\item Déterminer l'unique solution $f$ de ( $E$ ) telle que $f(0)=-1$.
	\end{enumMet}
	
\end{bmethodet}